% Generated draft: requires mathematical and source review.
\documentclass[11pt]{article}
\usepackage[margin=25mm]{geometry}
\usepackage{fontspec}
\setmainfont{Latin Modern Roman}
\usepackage{amsmath,amssymb,mathtools,unicode-math}
\setmathfont{Latin Modern Math}
\usepackage{enumitem,needspace,xurl,hyperref}
\setlength{\parindent}{0pt}
\setlength{\parskip}{5pt}
\setlength{\emergencystretch}{3em}
\newcommand{\N}{\mathbb{N}}
\newcommand{\Z}{\mathbb{Z}}
\newcommand{\Q}{\mathbb{Q}}
\newcommand{\R}{\mathbb{R}}
\newcommand{\C}{\mathbb{C}}
\newcommand{\F}{\mathbb{F}}
\newcommand{\A}{\mathbb{A}}
\newcommand{\PP}{\mathbb{P}}
\newcommand{\E}{\mathbb{E}}
\newcommand{\eps}{\varepsilon}
\newcommand{\dd}{\,\mathrm{d}}
\begin{document}
\section*{UnsolvedMath ID: 700441}
\textbf{Source problem code:} catalogue entry
\subsection{Definitions and mathematical statement}
Let \(\mathcal F\) denote the class of all finite groups. A group \(G\) is finitely presented if there exist nonnegative integers \(m,r\) and words \(w_1,\ldots,w_r\) in the free group \(F_m\) such that
\[
G\cong F_m/\langle\!\langle w_1,\ldots,w_r\rangle\!\rangle,
\]
where the denominator is the normal subgroup generated by the displayed words.

Write \(T_Q(\mathcal K)\) for the \(Q\)-theory of a class \(\mathcal K\) in the specific sense of Kourovka Notebook Problem 2.40. This is an unresolved source-defined term: the supplied record does not include Problem 2.40 or define \(Q\)-theory, so its logical fragment and semantics cannot be specified from the supplied evidence.

The question attributed to D. M. Smirnov is whether
\[
\exists\,G\quad\bigl(G\text{ is a finitely presented group}\ \land\ T_Q(\mathcal F)=T_Q(\{G\})\bigr).
\]
This preserves the requested equality as a statement schema; a fully specified mathematical formulation requires the definition referenced in Problem 2.40. No alternative meaning of \(Q\)-theory is assumed.
\subsection{Short English statement}
Does a single finitely presented group have exactly the same \(Q\)-theory as the class of all finite groups, using the meaning of \(Q\)-theory specified in Kourovka Notebook Problem 2.40? That definition is absent from the supplied record.
\subsection{Sources}
\begin{itemize}
\item Local UnsolvedMath registry, supplied complete source record: permanent ID \(700441\), exact source problem code \texttt{catalogue entry}, title \emph{Kourovka Notebook Problem 7.45}; upstream revision \texttt{d371ee917ac9647fbb06cd6fc529deec6f3dda82}. The record was created and updated on 2026-10-16 and labels the problem open; this is supplied registry metadata, not an independently verified mathematical or literature assessment.
\item \emph{The Kourovka Notebook (Unsolved Problems in Group Theory)}, Problem 7.45, proposed by D. M. Smirnov in the 7th issue (1980). The supplied record identifies the 21st issue edition (2026), arXiv:1401.0300v48. The collection description names E. I. Khukhro and V. D. Mazurov as editors.
\item The statement expressly refers to Kourovka Notebook Problem 2.40 for the meaning of \(Q\)-theory. Its text is not supplied, and neither that definition nor the cited edition has been independently consulted for this draft.
\end{itemize}
\end{document}
